\chapter{Why Modern Conceptual Tools Are Needed}
\label{ch:modern-tools}

\epigraph{\itshape A new scientific truth does not triumph by
convincing its opponents and making them see the light, but rather
because its opponents eventually die, and a new generation grows up
that is familiar with it.}{--- Max Planck, \textit{Scientific
Autobiography}, 1949}

\noindent
The previous three chapters have established the project. We have
seen (in \xref{ch:two-failures}) that contemporary Islamic economics
has failed to engage with the metaphysical claims of its tradition.
We have developed (in \xref{ch:levels-of-rigor}) the methodological
apparatus required to engage at varying degrees of rigor. We have
catalogued (in \xref{ch:textual-foundations}) the textual base from
which the claims arise.

The question now arises: why do we need \emph{modern} conceptual
tools? Why not simply continue the work the classical tradition
began? The Akbarian metaphysical system, after all, provides a
sophisticated articulation of the underlying ontology. The classical
\arb{kalam} tradition gave centuries of careful argument about
qadar and human action. What does the modern conceptual apparatus
add?

This chapter answers that question. The short answer is that the
classical tools, sophisticated as they are, were developed within a
framework that did not have access to certain features of reality
that twentieth-century mathematics and physics have made visible.
Quantum interpretations have shown that the classical Newtonian
picture of reality is incomplete, in ways that affect the questions
the classical theologians were asking. Ergodicity economics has
shown that standard expected-value reasoning, on which most economic
analysis rests, is structurally inadequate for processes with
absorbing barriers. Network science has shown that the linear
scaling intuitions on which conventional household economics rests
are violated by most real social systems. Information theory has
shown that the accumulation of nominal quantity is not the same as
the preservation of real value. Each of these is a development of
the last seventy-five years that the classical Islamic thinkers did
not have access to.

To use these tools is not to abandon the classical tradition. It is
to extend it with conceptual resources that the classical tradition
itself could have used had they been available. The modern thinker
who has both the classical apparatus and the modern apparatus is in
a better position to articulate the metaphysical claims than either
the classical thinker (who lacks the modern apparatus) or the modern
thinker who knows only the modern apparatus and does not see how
the classical tradition was already articulating these features.

The chapter proceeds in eight sections. \xref{sec:limits-neoclass}
explains the specific limitations of neoclassical economic
formalism. \xref{sec:complex-systems} discusses the rise of complex
systems theory as a methodological alternative.
\xref{sec:quantum-foundations-prelim} introduces the quantum
foundations material that Chapter 5 will develop in detail.
\xref{sec:ergodicity-prelim} introduces ergodicity economics, the
focus of Chapter 6. \xref{sec:idealism-return} discusses the return
of idealism in contemporary philosophy of mind, the focus of
Chapter 7. \xref{sec:negative-claim} states the negative claim that
unifies the modern apparatus: classical materialism with ergodic
averaging is not the only intellectually serious framework.
\xref{sec:not-concordism} forestalls the misreading that this
project is a form of concordism. \xref{sec:ahead-4} previews how
Part II will develop each of the conceptual tools.

\section{The limits of neoclassical formalism}
\label{sec:limits-neoclass}

The dominant formalism of contemporary economics --- neoclassical
microeconomics, with its rational agents, complete preferences,
expected utility maximization, and equilibrium analysis --- is the
language in which contemporary Islamic economics has been written.
The choice was reasonable in 1970; it is now a constraint. We need to
be specific about what the formalism cannot express.

\subsection{Preferences cannot be ontological}

In neoclassical theory, the consumer is described by a preference
relation $\succeq$ over a consumption set $\mathcal{X}$. The
preference is taken as a primitive: it is what the consumer happens
to want. The theory does not (and on its own terms cannot) ask why
the consumer wants what they want, where the preferences come from,
or whether there are preferences they ought to have but do not.

This is not a flaw in the theory. It is a deliberate methodological
restriction: the theory takes preferences as given because adjudicating
preferences would require commitments that the theory does not want
to make. The restriction is appropriate when the question being asked
is what choices the agent will make given the preferences. It is
inappropriate when the question is about the structure of preferences
themselves.

The Islamic tradition is making claims about the structure of
preferences. The claim that the believer ought to be grateful is not
a claim about the believer's actual preferences but about the
preferences they should have. The claim that the wealth obtained
through unlawful means produces less satisfaction than the wealth
obtained through lawful means is not a claim that the agent reveals
in their choices but a claim about the structure of the satisfaction
function itself. The neoclassical formalism cannot express these
claims because it treats the preference structure as primitive.

To express claims about the structure of preferences requires either
extending the formalism (with categories like
\emph{higher-order preferences}, \emph{normative preferences},
\emph{ethical structure}) or supplementing it with a framework that
gives ontological content to the preference structure. Both extensions
exist in the contemporary literature; neither is part of standard
neoclassical theory.

\subsection{Risk is treated by expected value}

In standard expected utility theory, an agent facing risk computes
the expected utility of available actions and chooses the action
with the highest expectation. The theory is mathematically elegant
and produces tractable analyses. It is also, as ergodicity economics
has shown, structurally inappropriate for processes with absorbing
barriers --- which is to say, for most processes that involve real
wealth.

The issue, treated in detail in Chapter 6, is that the expected value
of a multiplicative wealth process can grow indefinitely while the
\emph{time-averaged} growth rate experienced by any individual
participant is negative. The two quantities are different and the
expected value, which standard theory uses, is the wrong one for
agents who experience their wealth dynamics over time. The
correction is the use of the time-averaged growth rate, which
corresponds to the geometric mean of returns rather than the
arithmetic mean.

The Islamic tradition's prohibition of riba is, in the present
project's reading, an injunction against arrangements that exhibit
non-ergodic dynamics --- arrangements where the expected return
diverges from the actual experience. The neoclassical formalism, by
treating expected value as the relevant decision criterion, cannot
articulate the underlying critique. The critique requires the
ergodicity-economics extension, which we develop in Chapter 6.

\subsection{Rationality is computational}

Standard economic theory treats the rational agent as a constrained
optimizer. The agent has a utility function, a budget constraint,
and other constraints; rationality is the choice of the action that
maximizes utility subject to constraints. This conception of
rationality is, in computational terms, the execution of an
optimization algorithm.

The conception is not without difficulties. The Penrose argument
from G\"odel (treated in Chapter 8) suggests that human reasoning
includes elements that no computational procedure can capture. If
this is right, then the standard economic conception of rationality
misses an important feature of how humans actually reason about
choice. The Islamic tradition's emphasis on \arb{niyyah}
(intention) as constitutive of action is, on the present reading, a
recognition that what the agent does is not exhausted by the
computable optimization the standard theory describes. The treatment
in Chapter 14, of gratitude as constitutive of subsequent
flourishing, depends on this recognition.

\subsection{Time is treated additively}

Standard economic theory typically treats time additively: the
utility from a sequence of consumption is the discounted sum of the
period utilities. This is an analytical convenience that, again,
fits some processes well and others poorly. For multiplicative
wealth dynamics, it is structurally wrong: the appropriate
aggregation is the geometric mean (the time-averaged growth rate),
not the arithmetic sum.

The Islamic tradition's emphasis on the \emph{integral} of provision
over the lifetime --- the claim that no soul dies before completing
its allotted rizq --- is structurally a claim about time-averaged or
integrated quantities, not about per-period sums. The standard
formalism cannot express the claim cleanly; the integral formulation
required is what the ergodicity-economics framework provides.

\subsection{Causality is unidirectional}

Standard economic theory treats causal relations as unidirectional:
preferences cause choices; choices, given prices, cause demands;
aggregate demands, given supply, cause equilibrium prices. The
direction of causality is given by the model's structure.

The Akbarian tradition --- and the contemporary analytical idealism
that structurally parallels it --- treats causality as more
complicated. The agent's intention shapes the meaning of the action;
the action shapes the world; the world contains the structures within
which subsequent intentions form. The causal arrows go in both
directions and the structure is closed in a way that standard
economic models do not represent. The treatment of barakah in
Chapter 15 requires this richer notion of causality; the standard
formalism cannot express it.

\subsection{The aggregation of these limitations}

None of these limitations is fatal in isolation. Standard economic
theory has produced, and continues to produce, valuable results
within the scope of its applicability. The limitations matter for
the present project because the metaphysical claims of the Islamic
tradition lie precisely outside the scope of the standard formalism.
The claims concern:

\begin{itemize}[leftmargin=2em]
    \item The ontological structure of preferences (rather than
    their content).
    \item The dynamics of multiplicative wealth processes with
    absorbing barriers (rather than the static equilibrium of
    additive utilities).
    \item The constitutive role of intention (rather than the
    computational optimization of given preferences).
    \item The integral character of provision over a lifetime
    (rather than the per-period sum).
    \item The bidirectional structure of agent-world causality
    (rather than the unidirectional causal arrow of standard
    models).
\end{itemize}

To express the claims in their proper character, we need a richer
formalism. The richer formalism exists; it has been developed in
adjacent fields over the last several decades; the present project
borrows from it.

\section{The rise of complex systems theory}
\label{sec:complex-systems}

One of the major intellectual developments of the late twentieth
century is the rise of complex systems theory --- the study of
systems whose behavior at the aggregate level exhibits features (like
emergence, scaling laws, phase transitions, and non-ergodicity) that
are not present in or derivable from the behavior of the individual
components in isolation.

The field has its roots in nineteenth-century thermodynamics and
twentieth-century cybernetics, but its contemporary form emerged in
the 1980s through the work of the Santa Fe Institute and similar
research centers. Major contributors include Murray Gell-Mann (whose
later work on the foundations of probability is directly relevant to
the present project), Stuart Kauffman (whose work on self-organization
in evolutionary systems has parallels to the Akbarian conception of
\arb{tajalli}), Per Bak (whose self-organized criticality has
applications to financial systems), and many others.

The relevance of complex systems theory to the present project is
direct. The standard economic formalism assumes that aggregate
behavior is the sum of individual behaviors. Complex systems theory
has shown that this assumption is often false: aggregate behavior
exhibits scaling laws, non-equilibrium dynamics, and emergent
properties that the sum-of-individuals view cannot capture. The
Islamic tradition's claims about household and community dynamics
(particularly the children-bring-rizq claim treated in Chapter 11)
require the complex-systems framework to be properly articulated.

The specific tool from complex systems theory that we will use most
extensively is the \emph{scaling-law} analysis pioneered by Geoffrey
West, Luis Bettencourt, and others. The scaling-law approach has
shown that social systems (cities, organizations, markets) exhibit
scaling exponents that diverge systematically from the linear
($\beta = 1$) baseline. The exponents are characteristic of the
type of system: super-linear ($\beta > 1$) for innovation and
information-related quantities; sub-linear ($\beta < 1$) for
infrastructural quantities. The discovery of these exponents has
transformed the analysis of urban and social systems.

The application to households --- which is what we need for the
treatment of children-bring-rizq --- has been less developed in the
literature. The present project takes this application up
specifically and shows that Islamically-structured households
satisfy the conditions for super-linear scaling.

\section{Quantum foundations and the failure of classical materialism}
\label{sec:quantum-foundations-prelim}

The second major resource we draw on is the foundations of quantum
mechanics. This is the topic of Chapter 5, where we develop the
material in detail. The present preview indicates why the material
matters for the present project.

The foundational situation in quantum mechanics is the following.
The mathematical formalism --- Hilbert spaces, the Schr\"odinger
equation, the Born rule --- is empirically established to
extraordinary precision. The interpretation of the formalism --- what
the wave function represents, what happens at measurement, what
reality is like at the fundamental level --- is contested. The major
interpretations include:

\begin{itemize}[leftmargin=2em]
    \item The Copenhagen interpretation (Bohr, Heisenberg): the
    wave function is not a description of reality but a tool for
    predicting measurement outcomes; measurement collapse is
    fundamental and not further explicable.
    \item The Many-Worlds interpretation (Everett, Carroll): the
    wave function is real and never collapses; measurement is
    entanglement; all outcomes occur in different branches of a
    universal wave function.
    \item Bohmian mechanics (de Broglie, Bohm): particles have
    definite positions guided by a non-local pilot wave; the wave
    function is real; trajectories are deterministic.
    \item QBism (Fuchs, Mermin, Schack): the wave function
    represents an agent's beliefs; probability is irreducibly
    perspectival; different observers can have different valid wave
    functions for the same system.
    \item Objective collapse theories (Ghirardi-Rimini-Weber,
    Penrose): the Schr\"odinger equation is augmented with a
    spontaneous-collapse term; collapse is a real physical process.
\end{itemize}

The interpretive disputes are not minor. They differ on what
reality is, on what causation is, on what the role of the observer
is, and on what determinism means. They cannot all be true. They are
not adjudicable by experiment in the present state of physics; the
choice is at the level of philosophical commitment, informed by the
physics.

For the present project, three implications follow.

First, the dominant view among non-specialists --- that physics
has settled what reality is --- is wrong. The professional
disagreement is real and substantial. The reader who has been told
that ``science has answered'' the questions about reality is being
misled.

Second, the Bohmian and QBist interpretations specifically map onto
Islamic metaphysical claims. Bohmian mechanics gives us a coherent
formal framework in which trajectories are fully determined and
motion is real --- a structure that resolves the classical tension
between qadar and ikhtiyar. QBism gives us a perspectival framework
in which probability reflects agent belief, allowing divine knowledge
to be perfect while human knowledge remains finite, without
contradiction. Both interpretations were developed by physicists
without theological motivation; the structural match with Islamic
metaphysics is a discovery, not a construction.

Third, the conceptual room that the interpretive disputes open is
precisely the room that the metaphysical claims of the Islamic
tradition need. The classical materialist commitment that ``the
universe is just particles in motion governed by deterministic laws
and observers are irrelevant'' is not an established result of
physics. It is one philosophical position among several, and the
others are philosophically defensible.

The materialist commitment, when assumed, makes the Islamic
metaphysical claims look like incursions of theology into physics.
The materialist commitment, when seen as one option among several,
no longer has the standing to make those claims look anomalous.

\section{Ergodicity economics: the corrected formalism}
\label{sec:ergodicity-prelim}

The third major resource is ergodicity economics, the topic of
Chapter 6. The development was led by Ole Peters at the London
Mathematical Laboratory in collaboration with Murray Gell-Mann; the
foundational paper is Peters and Gell-Mann (2016).

The basic insight is that standard economic theory, descended from
Bernoulli (1738) and formalized by von Neumann and Morgenstern
(1944), uses the expected value (the ensemble average) as the
relevant decision criterion. This is appropriate for ergodic
processes (where the time average equals the ensemble average) and
inappropriate for non-ergodic processes (where they differ). Most
economic processes of interest are non-ergodic.

The applied consequences are significant. Standard finance theory
recommends maximizing expected return; this leads, with positive
probability, to ruin. The Kelly criterion (Kelly, 1956) provides the
correct decision rule for multiplicative wealth dynamics, recommending
the fraction of wealth that maximizes the geometric mean of returns
(equivalently, the time-averaged growth rate). The recommendations
differ substantially.

For the Islamic tradition's prohibition of riba, the ergodicity
analysis is decisive. Riba transactions, with default risk, are
multiplicative random processes with absorbing barriers. They are
non-ergodic. Their arithmetic expected return is positive (the
ensemble average grows), but the time average for any individual
participant is negative. The Quranic claim that Allah destroys riba
is, in literal mathematical sense, a true statement about the
process. The treatment in Chapter 13 develops this in detail.

Ergodicity economics also has positive implications. The rational
response to non-ergodic dynamics is risk pooling: the conversion of
a non-ergodic individual situation into an ergodic collective
situation through redistribution. This is structurally what
\arb{takaful} (Islamic mutual insurance) does, what \arb{zakat}
does at the macro scale, and what the broader package of Islamic
financial prohibitions amounts to. The ergodicity framework makes
visible the rational foundation of these institutions.

\section{The return of idealism in contemporary philosophy of mind}
\label{sec:idealism-return}

The fourth major resource is the contemporary defense of idealism
in philosophy of mind. The most rigorous current articulation is
the analytical idealism of Bernardo Kastrup, treated in detail in
Chapter 7.

Idealism, the metaphysical position that mind is fundamental and
matter is derivative, was the dominant position in Western
philosophy for centuries (in the work of Berkeley, Hegel, and many
others) and was largely abandoned in twentieth-century academic
philosophy in favor of materialism (or physicalism). The abandonment
was partly intellectual (the rise of the natural sciences seemed to
make materialism the default), partly sociological (philosophers
trained in materialist departments inherited the materialist
default), and partly philosophical (the so-called ``hard problem of
consciousness,'' Chalmers (1995), was not yet articulated as a
fatal challenge to materialism).

The hard problem of consciousness has, in the last three decades,
become an increasingly serious challenge to materialism. The
problem is that no description of brain processes appears to entail
the existence of subjective experience. The materialist holds that
brain processes give rise to experience; but no specification of
how, or why, has been forthcoming. This has led some philosophers to
propose that consciousness is a fundamental feature of matter
(panpsychism); others to propose that consciousness is fundamental
and matter is derivative (idealism). Kastrup's analytical idealism
is the most rigorous contemporary version of the latter.

The relevance for our project is direct. Islamic metaphysics is
fundamentally idealist. The doctrine of \arb{tawhid} asserts that
there is no being but Allah, that all created things are loci of
divine self-disclosure, and that the conventional materialist picture
of independently existing physical substances is, on the Akbarian
reading, a confusion. Until quite recently, this position was simply
out of step with academic philosophical respectability. Now it is
not. The contemporary defenders of idealism are not Muslims and have
no Islamic agenda; they are working within the mainstream tradition
of analytic philosophy. Their arrival at structurally similar
conclusions is independent confirmation that the Islamic position is
philosophically defensible.

The treatment in Chapter 7 develops the convergence in detail.

\section{The negative claim that unifies the modern apparatus}
\label{sec:negative-claim}

Each of the four resources just discussed --- complex systems
theory, quantum foundations, ergodicity economics, analytical
idealism --- is doing different work. They are not parts of a
single research program; they were developed in different fields by
different people for different reasons. What unites them, for the
purposes of the present project, is a shared \emph{negative claim}:

\begin{principlebox}[The unifying negative claim]
The dominant intellectual framework of late-twentieth-century
academic discourse --- classical materialism with ergodic averaging,
computational rationality, and observer-independent reality --- is
not the only intellectually serious framework. Each of its
component commitments is contested at the highest professional
level. The framework is one option among several; its dominance is
sociological, not epistemic.
\end{principlebox}

This negative claim is what the four resources, taken together,
establish. None of them, on its own, would suffice. Complex systems
theory, in isolation, could be read as just an extension of standard
economic methods. Quantum foundations, in isolation, could be read
as a debate confined to physics with no implications for economics.
Ergodicity economics, in isolation, could be read as a technical
correction to standard finance theory. Analytical idealism, in
isolation, could be read as one philosophical position among many.

Together, they establish something more substantial: that the
materialist-ergodic-computational framework is a constructed
intellectual edifice, not a forced conclusion from the data. The
framework can be coherently rejected. The Islamic tradition, working
from a different starting point, was already articulating the
features of reality that classical materialism had to ignore. With
the modern apparatus, the Islamic articulation can be expressed in
formal terms that the contemporary academic reader can engage with.

The negative claim is what makes the positive project possible. We
are not asking the contemporary academic reader to accept Islamic
theology; we are asking the reader to recognize that the framework
within which Islamic theology has been dismissed as scientifically
implausible is itself a contested framework, and that the
reformulation of Islamic claims in modern formal vocabulary is a
serious intellectual enterprise.

\section{This is not concordism}
\label{sec:not-concordism}

A persistent risk in projects of this kind is that they slide into
concordism. Concordism is the genre that attempts to demonstrate the
truth of a religious tradition by matching its texts against the
latest scientific findings. The genre has been particularly common
in the Islamic intellectual context over the last hundred years,
producing a literature that ``finds'' modern science predicted in
Quranic verses --- big bang cosmology in \Quran{21:30}, embryology
in \Quran{23:14}, plate tectonics in various verses, and so on.

Concordism is bad scholarship for several reasons. It reads modern
scientific concepts back into ancient texts, where they cannot have
been intended. It cherry-picks favorable scientific findings while
ignoring unfavorable ones. It makes the religious tradition hostage
to the scientific consensus of the moment, requiring revision
whenever the consensus shifts. It substitutes apologetics for
serious engagement.

The present project is structurally different from concordism in
several specific ways.

First, the project does not claim that the Quran or the hadith
\emph{predicted} modern findings. The convergences traced are not
predictions. They are structural recognitions: features of the
world that classical materialism failed to see, that modern science
has been forced by independent argument to acknowledge, and that the
Islamic tradition has been articulating in its own vocabulary because
it was working from different starting premises.

Second, the project does not depend on any particular contemporary
scientific consensus. The four resources we draw on are themselves
\emph{minority} positions within their fields. Bohmian mechanics is
not the consensus interpretation of quantum mechanics. Ergodicity
economics is not the consensus framework of mainstream finance.
Analytical idealism is not the consensus position in philosophy of
mind. The project draws on these minority positions because they are
the contemporary positions that map most clearly onto Islamic
metaphysical claims, not because they have been accepted by the
mainstream.

Third, the project is honest about the limits of the convergences.
The classical materialist framework does explain a great deal about
the world; the modern resources do not refute it across the board;
the Islamic metaphysical claims do not all admit of formalization;
some are Level 2 rather than Level 3; some remain open. Concordism,
by contrast, claims more than it can deliver and does not
acknowledge limits.

Fourth, the project is methodologically explicit. The three-level
taxonomy of \xref{ch:levels-of-rigor} is a discipline against
over-claiming. Each chapter of Part III flags the level of treatment
it is attempting and acknowledges what would be required for higher-level
claims. Concordism, by contrast, has no methodological discipline:
any convergence found is treated as confirmation, with no procedure
for distinguishing genuine convergence from coincidence.

The reader who suspects concordism is asked to evaluate the work on
these methodological criteria. If the criteria are met --- and the
author argues that they are --- then the work is doing something
different.

\section{Looking ahead: the structure of Part II}
\label{sec:ahead-4}

Part II of the book develops, in detail, each of the conceptual
resources surveyed in this chapter. The chapters of Part II are not
self-contained tutorials in the underlying fields --- they would be
much longer if they were --- but they are pedagogically careful
expositions of the specific material that Part III's worked
treatments will use.

\textbf{Chapter 5: Quantum Foundations and the Measurement Problem.}
The chapter develops the wave function, the Schr\"odinger equation,
the measurement problem, and the major interpretive disputes. The
treatment is sufficient for Chapter 16 (qadar) and provides the
philosophical context for the broader project.

\textbf{Chapter 6: Ergodicity Economics.} The chapter develops the
distinction between ensemble average and time average, formalizes
ergodicity, treats the Kelly criterion in detail, and applies the
framework to wealth dynamics. The treatment is the foundation for
Chapter 12 (halal vs.\ haram dynamics) and Chapter 13 (riba as
financial entropy).

\textbf{Chapter 7: Analytical Idealism and the Ontology of Mind.}
The chapter treats the hard problem of consciousness, the
contemporary positions in philosophy of mind, and Kastrup's
analytical idealism in detail. The treatment provides the
contemporary philosophical register that Chapter 9's treatment of
Akbarian metaphysics will compare itself to.

\textbf{Chapter 8: Non-Computability and the Reality of Intention.}
The chapter develops G\"odel's incompleteness theorems, Turing's
halting problem, the Penrose argument from G\"odel, and the
Penrose-Hameroff Orch-OR theory. The treatment is the foundation
for the discussions of niyyah throughout Part III, particularly in
Chapter 14.

\textbf{Chapter 9: Akbarian Metaphysics and the Self-Disclosure of
Wujud.} The chapter develops the Akbarian system in the depth required
for the worked treatments. Wujud, the divine names, tajalli, the
imaginal world, the perfect human --- each is treated as a structured
metaphysical doctrine and connected to the contemporary frameworks
introduced in the preceding chapters.

The reader should expect Part II to be more technical than Part I.
The technical material is unavoidable: rigorous treatment of the
metaphysical claims requires the technical apparatus, and the
apparatus must be developed before it can be used. Readers with
prior background in any of the relevant fields can read the
corresponding chapter quickly; readers without can take their time.

After Part II, the worked treatments of Part III will draw freely on
the apparatus, with cross-references back to the relevant chapters
of Part II. The reader who keeps the structure in mind --- Part I:
foundations; Part II: tools; Part III: applications; Part IV:
synthesis --- will find the cumulative argument easier to follow.

\section{Conclusion of Part I}

Part I of the book is now complete. The reader has been shown:

\begin{enumerate}[leftmargin=2em]
    \item Why the present work is necessary (Chapter 1).
    \item The methodological discipline that organizes it (Chapter
    2).
    \item The textual base from which the metaphysical claims arise
    (Chapter 3).
    \item Why modern conceptual tools are needed and which tools
    will be developed (Chapter 4).
\end{enumerate}

The reader is now prepared to engage with Part II, where the
conceptual apparatus is developed in detail. The careful reader will
find that the apparatus, once mastered, opens the metaphysical
claims of the Islamic tradition to forms of analysis they have not
previously received.

A note on what follows in subsequent installments. The present
volume contains the front matter and Part I. Part II (Chapters 5
through 9) will be the next installment; it develops the modern
conceptual apparatus. Part III (Chapters 10 through 17) will follow;
it treats the eight canonical metaphysical claims with the apparatus
developed in Part II. Part IV (Chapters 18 through 21) will conclude
the volume; it synthesizes the worked treatments and lays out the
research agenda that the synthesis implies.

The author takes the present moment to thank the reader who has
followed Part I to its conclusion. The introductory work has been
necessary --- without it, the substantive treatments to come would
be unintelligible --- but it has also been demanding. The chapters
that follow are, in the author's view, more rewarding intellectually:
they take the apparatus established in Part I and use it to do
genuinely new work. The reader who has stayed with the project
through Part I is now ready for that work.

\begin{summarybox}[Chapter summary]
\textbf{The neoclassical formalism cannot express}: the ontological
structure of preferences, the dynamics of multiplicative wealth
processes with absorbing barriers, the constitutive role of
intention, the integral character of provision over a lifetime, or
the bidirectional structure of agent-world causality.

\textbf{Modern conceptual resources to be developed in Part II}:
\begin{itemize}[leftmargin=1.5em]
    \item Complex systems theory and scaling laws (background for
    Chapter 11).
    \item Quantum foundations and interpretive disputes (Chapter 5,
    used in Chapter 16).
    \item Ergodicity economics (Chapter 6, used in Chapters 12 and
    13).
    \item Analytical idealism (Chapter 7, compared with Chapter 9).
    \item Non-computability and the Penrose argument (Chapter 8,
    used in Chapter 14).
\end{itemize}

\textbf{The unifying negative claim}: classical materialism with
ergodic averaging is not the only intellectually serious framework;
each of its component commitments is contested at the highest
professional level.

\textbf{Why this is not concordism}: the project does not claim that
the Quran predicted modern findings; the convergences are structural
recognitions, not predictions; the project draws on minority
positions in contemporary fields rather than mainstream consensus;
the project is methodologically explicit and acknowledges limits.

\textbf{What comes next}: Part II (Chapters 5--9) will develop the
modern conceptual apparatus in detail.
\end{summarybox}

\cleardoublepage
